% Chapter 3 — Regret-Minimizing Algorithms

\chapter{Regret-Minimizing Algorithms}
\label{chap:algorithms}

This chapter introduces the algorithms used in the experiments. All methods are
taken from the literature; the purpose here is only to state the implemented
variant, the main idea behind it, and the theoretical guarantee that is most
relevant for the later empirical comparison. Detailed proofs are not repeated.

The implementations fall into three groups: standalone external-regret
learners, regret-matching procedures, and reductions that construct stronger
regret guarantees from external-regret sublearners.

% ============================================================
\section{External-Regret Learners}

\paragraph{Hedge and Optimistic Hedge.}

Hedge is the main full-information external-regret baseline. The implementation samples according
to exponential weights and, for a known horizon, uses
\(\eta=\sqrt{\frac{8\log K}{T}}\) \citep{freund_decision-theoretic_1997}. With bounded rewards this gives the standard
root-$T$ distribution-regret guarantee. The same fixed-horizon Hedge learner is used inside BM-Hedge.
Ito-Hedge instead uses the anytime version
\(\eta_s=\sqrt{8\log K/s}\), where \(s\) is the local update count of the
selected Ito sublearner.

Following \citet{chen_hedging_2020}, Optimistic Hedge (OptHedge) uses the optimistic update
\[
    p_{t+1}(i)\propto
    p_t(i)\exp\!\left(\eta(2r_t(i)-r_{t-1}(i))\right)
\]
with learning rate \(\eta=\left(\frac{\log K}{T}\right)^{1/6}\). Their fast
\(O(T^{1/6}\log^{5/6}K)\) mixed external-regret result applies when both
players use Optimistic Hedge in a two-player repeated game. It is not a
general adversarial guarantee, so the one-player OptHedge experiments are used
only as empirical comparisons.

\paragraph{EXP3, EXP3-IX, and Tsallis-INF.}

EXP3 is the explicit-exploration bandit baseline. Following
\citet{auer_nonstochastic_2002}, the standalone implementation uses
\[
    \gamma=\min\left\{1,\sqrt{\frac{K\log K}{(e-1)T}}\right\},
    \qquad
    \eta=\frac{\gamma}{K},
\]
together with the usual importance-weighted reward estimate. The EXP3 learners
inside BM-EXP3 use the same update rule but the Blum--Mansour tuning
\(\gamma=\min\{1,\sqrt{K\log K/T}\}\) \citep{blum_external_2007}.

EXP3-IX replaces explicit uniform exploration by implicit exploration. In loss form, the estimator divides the observed
loss by \(p_t(i)+\gamma\), and no additional uniform mixture is used. Following
\citet{neu_explore_2015}, the implementation uses
\[
    \eta=\sqrt{\frac{2\log K}{KT}},
    \qquad
    \gamma=\frac{\eta}{2}.
\]
Neu proves a root-$T$-order high-probability realized external-regret bound.

Tsallis-INF is used only as the inner bandit learner of Ito-Tsallis. Following
\citet{zimmert_tsallis-inf_2021}, the implementation uses the
importance-weighted \(\alpha=1/2\) variant with local learning rate
\(\eta_s=1/\sqrt{s}\). This is an equivalent scaling of their formulation, which gives adversarial pseudo-regret of
root-$T$ order.

% ============================================================
\section{Regret Matching}

Both regret-matching implementations maintain cumulative pairwise replacement
gains
\[
    C_t(i,j)
    =
    \sum_{s=1}^{t}
    \mathbf 1\{I_s=i\}
    \bigl(r_s(j)-r_s(i)\bigr).
\]
With rewards in \([0,1]\), the implementation sets the fixed normalization to \(K\),
which is a valid choice for the Hart--Mas-Colell procedure \citep{hart_simple_2000},
and forms a stochastic transition matrix using
\[
    Q_t(i,j)=\frac{[C_t(i,j)]_+}{Kt},
    \qquad i\neq j,
\]
with diagonal entries chosen so that each row sums to one.

Regret Matching (RM) implements the adaptive procedure of Section~2 of
\citet{hart_simple_2000}: after playing action \(i\), the next distribution is
the \(i\)-th row of \(Q_t\). Hart and Mas-Colell prove that if all players use
this procedure, the empirical distribution of play converges almost surely to
the correlated-equilibrium set. This is a joint self-play result, not a
general individual no-internal-regret guarantee against arbitrary adaptive
opponents.

Stationary Regret Matching (SRM) instead plays an invariant distribution of
the transition matrix,
\[
    p_{t+1}=p_{t+1}Q_t.
\]
Theorem~A of \citet{hart_simple_2000} gives vanishing individual pairwise
regret almost surely under the paper's conditional-independence assumption.
This distinction between RM and SRM is important in the later one-player and
cross-play comparisons.

% ============================================================
\section{Swap-Regret Reductions}

\paragraph{Blum--Mansour.}

The Blum--Mansour reduction maintains \(K\) external-regret learners, one for
each source action \citep{blum_external_2007}. Their output distributions form
a stochastic matrix \(Q_t\), and the master distribution is chosen stationary:
\[
    p_t=p_tQ_t.
\]
Each row can be interpreted as learning a replacement distribution for one
source action.

In the full-information version, all rows are updated every round with the
weighted reward vector \(p_t(i)r_t\). BM-Hedge uses the fixed-horizon Hedge
learner above. BM-OptHedge replaces Hedge by Optimistic Hedge and, following
\citet{chen_hedging_2020}, uses
\[
    \eta=
    \left(\frac{K\log K}{m^2T}\right)^{1/4}.
\]
In repeated games, \(m\) is the actual number of players; the one-player
benchmark keeps the implementation default \(m=2\). The fast
\(O((K\log K)^{3/4}m^{1/2}T^{1/4})\) mixed swap-regret guarantee of
\citet{chen_hedging_2020} applies when all players use this optimistic
self-play dynamic.

BM-EXP3 is the bandit Blum--Mansour variant. It uses EXP3 row learners and
updates all \(K\) rows every round using the importance-weighted feedback
specified by \citet{blum_external_2007}. The corresponding guarantee is on
swap pseudo-regret.

\paragraph{Ito reduction.}

Ito retains the stationary-distribution construction but updates only one row
per round \citep{ito_tight_2020}. First
\[
    J_t\sim p_t,
\]
then the played action is sampled from row \(J_t\), and only that selected
sublearner receives feedback. Because \(p_t=p_tQ_t\), the marginal
distribution of the played action is still \(p_t\).

Ito-Hedge uses the anytime Hedge learner, while Ito-Tsallis uses
importance-weighted Tsallis-INF. Their learning rates depend on the selected
row's local update count rather than on the global round. Ito's generic
reduction shows that an external learner with expected regret \(r(S)\) yields
\(K\,r(T/K)\) for the maximization-outside-expectation swap criterion
\citep{ito_tight_2020}. In the terminology of
Section~\ref{sec:regret-notions}, this is swap pseudo-regret.

\paragraph{LCE-IX.}

LCE-IX is the bandit swap-regret method of
\citet{huang_near-optimal_2023}. Like Blum--Mansour it uses a stationary
distribution and updates all rows every round, but combines this construction
with implicit exploration. Following the parameter schedule in the same work, the implementation uses
\[
    \eta_t=\sqrt{\frac{\log K}{t}},
    \qquad
    \gamma_t=\frac{\eta_t}{2}.
\]
Huang and Pan prove \(O(K\sqrt{T\log K})\) expected swap regret---an expected
action-regret quantity in the terminology of Chapter~\ref{chap:background}---
and a corresponding high-probability bound for instantaneous swap regret.

% ============================================================
\section{Algorithm Summary}
\label{sec:algorithm-summary}

Table~\ref{tab:algorithm-summary} summarizes the role of each implementation
and the literature guarantee most closely associated with the implemented
parameterization. All displayed big-$O$ rates are cumulative in the horizon $T$ and
suppress constant factors; the RM and SRM entries instead state their
asymptotic average-regret guarantees. They are not directly comparable unless the regret
object and the assumptions indicated below the table agree.

\begin{table}[htbp]
    \centering
    \footnotesize
    \caption{Algorithms used in the experiments and their principal literature guarantees.}
    \label{tab:algorithm-summary}
    \begin{tabular}{@{}llp{2.55cm}p{3.15cm}p{3.65cm}@{}}
        \toprule
        Algorithm & Feedback & Main role & Regret notion & Literature guarantee \\
        \midrule
        Hedge
        & Full
        & External baseline
        & External distribution regret
        & $O(\sqrt{T\log K})$ \\
        OptHedge\textsuperscript{*}
        & Full
        & Optimistic external regret
        & External distribution regret
        & $O(T^{1/6}\log^{5/6}K)$ \\
        RM\textsuperscript{**}
        & Full
        & Adaptive regret matching
        & Internal action regret
        & $R_T^{\mathrm{int}}/T\to0$ a.s. \\
        SRM\textsuperscript{**}
        & Full
        & Stationary pairwise regret
        & Internal action regret
        & $R_T^{\mathrm{int}}/T\to0$ a.s. \\
        BM-Hedge
        & Full
        & BM swap reduction
        & Swap distribution regret
        & $O(\sqrt{TK\log K})$ \\
        BM-OptHedge\textsuperscript{*}
        & Full
        & Optimistic BM reduction
        & Swap distribution regret
        & $O\!\left(m^{1/2}(K\log K)^{3/4}T^{1/4}\right)$ \\
        Ito-Hedge
        & Full
        & One-row Ito reduction
        & Swap pseudo-regret
        & $O(\sqrt{TK\log K})$ \\
        \midrule
        EXP3\textsuperscript{***}
        & Bandit
        & External baseline
        & External pseudo-regret
        & $O(\sqrt{TK\log K})$ \\
        EXP3-IX\textsuperscript{****}
        & Bandit
        & Implicit-exploration baseline
        & Realized external action regret
        & $O(\sqrt{TK\log K})$ w.h.p. \\
        BM-EXP3
        & Bandit
        & BM swap reduction
        & Swap pseudo-regret
        & $O(K\sqrt{KT\log K})$ \\
        Ito-Tsallis
        & Bandit
        & One-row Ito reduction
        & Swap pseudo-regret
        & $O(K\sqrt{T})$ \\
        LCE-IX
        & Bandit
        & Stationary IX swap regret
        & Expected swap action regret
        & $O(K\sqrt{T\log K})$ \\
        \bottomrule
    \end{tabular}

    \vspace{0.5em}
    \begin{minipage}{0.97\textwidth}
        \footnotesize
        \textsuperscript{*}The optimistic fast rates are repeated-game results rather than generic adversarial guarantees: the OptHedge bound assumes two-player self-play with both players using Optimistic Hedge, while the BM-OptHedge bound assumes all $m$ players use the optimistic BM dynamic \citep{chen_hedging_2020}.\\
        \textsuperscript{**}For RM, the Hart--Mas-Colell Main Theorem assumes that all players use the adaptive procedure. For SRM, the individual vanishing-regret theorem allows arbitrary opponent behavior subject to conditional independence of players' current randomizations given the past \citep{hart_simple_2000}.\\
        \textsuperscript{***}The displayed EXP3 rate is the pseudo-regret guarantee for the nonstochastic adversarial-bandit model of \citet{auer_nonstochastic_2002}.\\
        \textsuperscript{****}For EXP3-IX the displayed order suppresses the confidence-dependent term in \citet{neu_explore_2015}; the theorem is a high-probability bound on realized external regret.
    \end{minipage}
\end{table}

The central implementation differences for the empirical comparison are
therefore simple: standalone learners target external regret; RM and SRM react
directly to pairwise replacement regret; BM and LCE-IX update all row learners
through a stationary construction; and Ito updates only one selected row per
round. The next chapter describes how these methods are evaluated under a
common experimental framework.
